\begin{tablalafrox}{Diccionario del Producto 10: Servicio de operación y salida. Fuente: según cita por cuenta en cada encabezado.}{tab:T14-dic-prod-10}
  {L{1.2cm} Y Y Y C{0.9cm} L{2.3cm}}
  {\cab{Código} & \cab{Paquete de trabajo} & \cab{Entregable} & \cab{Criterio de aceptación} & \cab{Resp.} & \cab{Período}}
  \multicolumn{6}{>{\columncolor{lafroxBg}\RaggedRight\arraybackslash}p{\dimexpr\textwidth-2\tabcolsep\relax}}{\textbf{Cuenta 10.1 — Servicios de operación.} Esta cuenta reúne los servicios que mantienen la solución funcionando cada día. {\color{lafroxGris}Fuente: (Bases Administrativas, art. 17, pp. 12–13, y Formulario T-14, p. 63); (Bases Técnicas del caso, secciones 17.2–17.6, pp. 32–34).}} \tabularnewline*
  10.1.1 & Servicio de centro de operaciones (NOC 24\ensuremath{\times}7) y mesa de servicio & El servicio prestado durante los 36 meses, con un informe mensual de servicio y un informe anual por cada año de operación. & Se cumplen los horarios de cobertura (NOC 24\ensuremath{\times}7 y mesa de 04:00 a 22:00, con 24\ensuremath{\times}7 en septiembre y diciembre) en cada uno de los 36 meses. & SRE & Operación · Meses 21 a 56. Se paga en el hito mensual. \\
  10.1.2 & Servicio de gestión de incidentes, problemas y niveles de servicio (Art. 78°) & El servicio prestado durante los 36 meses, con un informe mensual de niveles de servicio, los informes de causa raíz y un informe anual. & Se cumplen los tiempos, y cada incidente crítico tiene su informe a tiempo. & SRE & Operación · Meses 21 a 56. Las multas se descuentan del hito mensual. \\
  10.1.3 & Servicio de continuidad y pruebas de recuperación ante desastres & Siete pruebas de conmutación en los 36 meses (una al inicio de la Operación y luego dos por año, la última en el mes 54), cada una con su informe. & Cada prueba alcanza el RTO y el RPO comprometidos en una conmutación real, programada fuera de los congelamientos y con separación máxima de seis meses. & SRE & Operación · Meses 21, 27, 33, 39, 45, 51 y 54. \\
  10.1.4 & Servicio de centro de operaciones de seguridad (SOC 24\ensuremath{\times}7) con SIEM & El servicio prestado desde el mes 13 hasta el mes 56, con un informe mensual de seguridad. & Cobertura 24\ensuremath{\times}7 todos los meses, cada incidente crítico avisado dentro de dos horas y cada brecha notificada dentro de 24 horas. & SEG & Operación · Meses 13 a 56. \\
  10.1.5 & Informes mensuales de restauración de respaldos & 36 informes de restauración, uno por mes. & Cada mes la restauración funciona sin errores, con su tiempo medido. & SRE & Operación · Mensual, meses 21 a 56. \\
  \midrule
  \multicolumn{6}{>{\columncolor{lafroxBg}\RaggedRight\arraybackslash}p{\dimexpr\textwidth-2\tabcolsep\relax}}{\textbf{Cuenta 10.2 — Mantenimiento.} Esta cuenta de control reúne el trabajo que mantiene la solución sana y actualizada, separado por tipo: software, ciberseguridad y hardware. {\color{lafroxGris}Fuente: (Bases Administrativas, art. 17, pp. 12–13, y Formulario T-14, p. 63); (Bases Técnicas del caso, secciones 17.2–17.6, pp. 32–34).}} \tabularnewline*
  10.2.1 & Servicio de mantenimiento correctivo y evolutivo de software & Versiones desplegadas con cargo a la bolsa anual de horas, cada una con su registro: requisito, cambio, prueba, aprobación y despliegue. & Ningún despliegue cae en una fecha prohibida, y cada cambio tiene su registro completo. & DES & Operación · Meses 21 a 56. \\
  10.2.2 & Servicio de mantenimiento de ciberseguridad y resiliencia & En los 36 meses, los informes periódicos de vulnerabilidades, tres informes de intrusión (uno por año) y seis de resiliencia (dos por año). & Se cumplen los plazos de corrección; la prueba anual termina sin hallazgos críticos ni altos abiertos; hay dos pruebas de resiliencia con inyección de fallas por año, con separación máxima de seis meses. & SEG & Operación · Intrusión en meses 28, 40 y 52; resiliencia en meses 26, 32, 38, 44, 50 y 55. \\
  10.2.3 & Servicio de mantenimiento preventivo y correctivo de hardware & Registro de mantenimientos de los 36 meses (sala técnica, racks, gabinetes de borde y equipamiento de terreno) y un informe anual por cada año de operación. & Se cumplen los mantenimientos programados, y ningún equipo crítico queda sin repuesto. & SRE & Operación · Meses 21 a 56. \\
  10.2.4 & Registro e informes trimestrales de reducción de deuda técnica & Registro de deuda técnica e informe trimestral de su reducción. & El registro está cuantificado y se cumple la fracción de capacidad declarada. & DES & Operación · Meses 21 a 56. \\
  10.2.5 & Actualizaciones anuales de componentes de base & Plan anual de actualización y su ejecución. & Cada año hay un plan aprobado, ejecutado en su ventana y con su reversión ensayada. & SRE & Operación · Una vez por año, meses 21 a 56. \\
  10.2.6 & Informes semestrales de revisión eléctrica de la sala técnica & Seis informes de revisión eléctrica. & Cada semestre hay una revisión con su informe y sus correcciones. & SRE & Operación · Semestral, meses 21 a 56. \\
  10.2.7 & Plan y registro de ciclo de vida del equipamiento & Plan de ciclo de vida y registro de reposiciones y disposiciones. & Cada tipo de equipo del T-11 tiene su ciclo de vida, y cada equipo retirado tiene su disposición registrada. & SRE & Operación · Meses 21 a 56. \\
  \midrule
  \multicolumn{6}{>{\columncolor{lafroxBg}\RaggedRight\arraybackslash}p{\dimexpr\textwidth-2\tabcolsep\relax}}{\textbf{Cuenta 10.3 — Informes y obligaciones periódicas.} Esta cuenta de control reúne los informes y las tareas que las Bases fijan con una periodicidad durante la Operación y que no forman parte de los servicios diarios. {\color{lafroxGris}Fuente: (Bases Administrativas, art. 17, pp. 12–13, y Formulario T-14, p. 63); (Bases Técnicas del caso, secciones 17.2–17.6, pp. 32–34).}} \tabularnewline*
  10.3.1 & Informes trimestrales de gestión y proyección de capacidad & 12 informes trimestrales de capacidad. & Cada trimestre hay una proyección, con alertas y una propuesta de ajuste cuando corresponde. & SRE & Operación · Trimestral, meses 21 a 56. \\
  10.3.2 & Informes mensuales de consumo de nube (FinOps) & 36 informes mensuales de consumo. & Cada informe desglosa el consumo, y toda desviación sobre el presupuesto tiene su alerta y su explicación. & SRE & Operación · Mensual, meses 21 a 56. \\
  10.3.3 & Informes anuales de huella de carbono de la operación & Tres informes anuales de huella de carbono. & Cada informe declara su metodología y compara con el año anterior. & SRE & Operación · Anual, meses 21 a 56. \\
  10.3.4 & Actualizaciones semestrales del depósito de custodia de fuentes & Seis actualizaciones del depósito, con su constancia. & Cada semestre el depósito coincide con la versión en producción. & JP & Operación · Semestral, meses 21 a 56. \\
  \midrule
  \multicolumn{6}{>{\columncolor{lafroxBg}\RaggedRight\arraybackslash}p{\dimexpr\textwidth-2\tabcolsep\relax}}{\textbf{Cuenta 10.4 — Capacitación y transferencia en operación.} Esta cuenta de control reúne la capacitación y el traspaso de conocimiento que las Bases exigen después de la implementación. {\color{lafroxGris}Fuente: (Bases Administrativas, art. 17, pp. 12–13, y Formulario T-14, p. 63); (Bases Técnicas del caso, secciones 17.2–17.6, pp. 32–34).}} \tabularnewline*
  10.4.1 & Jornadas anuales de actualización y capacitación de personal nuevo & Seis jornadas de actualización como mínimo y registro de la capacitación del personal nuevo. & Cada año hay al menos dos jornadas, y cada persona nueva está capacitada en su perfil. & IMP & Operación · Meses 21 a 56. \\
  10.4.2 & Programa de mentoría al equipo técnico del CLIENTE & Plan de mentoría y registro de las sesiones. & La mentoría cubre al menos seis meses, y el equipo de TI ejecuta solo sus funciones al final del período. & IMP & Operación · Meses 21 a 27. \\
  10.4.3 & Base de conocimiento actualizada e informes de uso & Base de conocimiento actualizada e informe trimestral de uso. & Cada incidente resuelto con una solución nueva tiene su entrada, y hay métricas de uso y de utilidad cada trimestre. & SRE & Operación · Meses 21 a 56. \\
  \midrule
  \multicolumn{6}{>{\columncolor{lafroxBg}\RaggedRight\arraybackslash}p{\dimexpr\textwidth-2\tabcolsep\relax}}{\textbf{Cuenta 10.5 — Salida al término del contrato.} Esta cuenta de control reúne el traspaso final y la eliminación de los datos personales. {\color{lafroxGris}Fuente: (Bases Administrativas, art. 17, pp. 12–13, y Formulario T-14, p. 63); (Bases Técnicas del caso, secciones 17.2–17.6, pp. 32–34).}} \tabularnewline*
  10.5.1 & Traspaso de reversibilidad y acompañamiento de salida & LafroX ejecuta el Plan de Reversibilidad y documenta el traspaso en un acta. & El receptor opera sin interrupción durante los 90 días de acompañamiento y firma la recepción. & JP & Operación · Meses 54 a 56: el acompañamiento empieza al menos 90 días antes del término del contrato, mientras LafroX mantiene la operación. Si el CLIENTE fija un término anticipado o una transición posterior, el período se ajusta para conservar los 90 días. \\
  10.5.2 & Certificado de devolución o eliminación de datos personales & Certificado de devolución, transferencia de custodia y eliminación. & Los datos se devuelven íntegros en formatos abiertos; las copias en ambientes y respaldos de LafroX se eliminan; los objetos con retención inmutable (Object Lock, registro de auditoría de siete años) se transfieren en custodia al CLIENTE sin alterar su retención, y su eliminación se certifica al vencer cada retención. El certificado identifica qué se eliminó, qué se transfirió y con qué fecha de vencimiento. & SEG & Operación · Al término del contrato, después de 10.5.1. \\
\end{tablalafrox}
